\section{Bourbaki's locally convex extension of Banach--Alaoglu}

Throughout this chapter $E$ is a real locally convex topological vector space:
an $\mathbb{R}$-vector space with a topology making addition and scalar
multiplication continuous and admitting a base of convex neighbourhoods of the
origin. We write $\mathrm{WeakDual}\,\mathbb{R}\,E$ for the space of continuous
linear functionals $\varphi : E \to \mathbb{R}$ equipped with the weak-star
topology, i.e.\ the topology of pointwise convergence induced by the evaluation
maps $\varphi \mapsto \varphi(x)$.

\begin{definition}[Weak-star polar]
    \label{def:weak-star-polar}
    Let $s \subseteq E$. The \emph{weak-star polar} of $s$ is the set
    \[
        \operatorname{polar}(s)
        = \{\, \varphi \in \mathrm{WeakDual}\,\mathbb{R}\,E \mid
              \forall x \in s,\ \lVert \varphi(x) \rVert \le 1 \,\}.
    \]
    (This matches the Lean definition \texttt{LeanEval.Analysis.weakStarPolar}.)
\end{definition}

The proof realises each functional by its family of values and transfers the
problem to the space $E \to \mathbb{R}$ carrying the topology of pointwise
convergence. The first lemma records that this transfer is faithful.

\begin{lemma}[The evaluation embedding]
    \label{lem:coe-embedding}
    The coercion
    $(\cdot) : \mathrm{WeakDual}\,\mathbb{R}\,E \to (E \to \mathbb{R})$,
    $\varphi \mapsto (x \mapsto \varphi(x))$, is a topological embedding.
    Consequently, for any $S \subseteq \mathrm{WeakDual}\,\mathbb{R}\,E$, the set
    $S$ is compact if and only if its image $(\cdot)\,''\,S$ is compact in
    $E \to \mathbb{R}$.
\end{lemma}
\begin{proof}
    By definition the weak-star topology on $\mathrm{WeakDual}\,\mathbb{R}\,E$
    is the topology induced by this coercion, and the coercion is injective
    because a linear functional is determined by its values. An injective map
    from a space carrying the induced topology is a topological embedding, and a
    set is compact if and only if its image under a topological embedding is
    compact.
\end{proof}

\begin{lemma}[Pointwise boundedness via absorbency]
    \label{lem:polar-pointwise-bounded}
    \uses{def:weak-star-polar}
    Let $U$ be a neighbourhood of $0$ in $E$ and let $x \in E$. Then there
    exists $r \ge 0$ such that $\lVert \varphi(x) \rVert \le r$ for every
    $\varphi \in \operatorname{polar}(U)$.
\end{lemma}
\begin{proof}
    A neighbourhood of the origin in a topological vector space is absorbent, so
    there is a scalar $c \ne 0$ with $c^{-1} \cdot x \in U$. For any
    $\varphi \in \operatorname{polar}(U)$, linearity gives
    $\varphi(x) = c \, \varphi(c^{-1} \cdot x)$, and
    $\lVert \varphi(c^{-1} \cdot x) \rVert \le 1$ since $c^{-1} \cdot x \in U$.
    Hence $\lVert \varphi(x) \rVert \le \lvert c \rvert$, and $r = \lvert c
    \rvert$ works.
\end{proof}

\begin{lemma}[Bounded linear functionals are continuous]
    \label{lem:linear-bounded-continuous}
    Let $U$ be a neighbourhood of $0$ in $E$ and let $g : E \to \mathbb{R}$ be a
    linear functional with $\lVert g(x) \rVert \le 1$ for all $x \in U$. Then
    $g$ is continuous; equivalently, $g$ is the coercion of an element of
    $\mathrm{WeakDual}\,\mathbb{R}\,E$.
\end{lemma}
\begin{proof}
    The functional $g$ maps the neighbourhood $U$ of $0$ into the interval
    $[-1,1] \subseteq \mathbb{R}$. This interval is metric-bounded by
    \texttt{Metric.isBounded\_Icc}, hence von Neumann bounded by
    \texttt{NormedSpace.isVonNBounded\_of\_isBounded}. A linear map sending some
    neighbourhood of $0$ into a von Neumann bounded set is continuous, so $g$ is
    continuous and therefore is the coercion of a continuous linear functional,
    i.e.\ an element of $\mathrm{WeakDual}\,\mathbb{R}\,E$.
\end{proof}

\begin{lemma}[The image as an intersection]
    \label{lem:polar-image-eq}
    \uses{def:weak-star-polar, lem:linear-bounded-continuous}
    Let $U$ be a neighbourhood of $0$ in $E$. Write
    $L \subseteq (E \to \mathbb{R})$ for the range of the coercion
    $\mathrm{Hom}_{\mathbb{R}}(E, \mathbb{R}) \to (E \to \mathbb{R})$ that forgets
    linearity, i.e.\ the set of $\mathbb{R}$-linear functions $E \to \mathbb{R}$.
    Then
    \[
        (\cdot)\,''\,\operatorname{polar}(U)
        = L \cap \{\, g : E \to \mathbb{R} \mid
              \forall x \in U,\ \lVert g(x) \rVert \le 1 \,\}.
    \]
\end{lemma}
\begin{proof}
    Forward inclusion: if $g = (\cdot)\,\varphi$ for some
    $\varphi \in \operatorname{polar}(U)$, then $g$ is linear (as the coercion of
    a continuous linear functional) and $\lVert g(x) \rVert \le 1$ for all
    $x \in U$ by definition of $\operatorname{polar}(U)$. Backward inclusion:
    given a linear $g$ with $\lVert g(x) \rVert \le 1$ for all $x \in U$,
    Lemma~\ref{lem:linear-bounded-continuous} shows $g$ is the coercion of some
    $\varphi \in \mathrm{WeakDual}\,\mathbb{R}\,E$, and this $\varphi$ lies in
    $\operatorname{polar}(U)$ because $\lVert \varphi(x) \rVert = \lVert g(x)
    \rVert \le 1$ for $x \in U$; hence $g$ is in the image.
\end{proof}

\begin{lemma}[Closedness of the intersection]
    \label{lem:polar-image-closed}
    The set
    \[
        L \cap \{\, g : E \to \mathbb{R} \mid
              \forall x \in U,\ \lVert g(x) \rVert \le 1 \,\}
    \]
    from Lemma~\ref{lem:polar-image-eq} is closed in $E \to \mathbb{R}$ for the
    topology of pointwise convergence.
\end{lemma}
\begin{proof}
    The set $L$ of linear functions is closed by
    \texttt{LinearMap.isClosed\_range\_coe}. The second factor is the
    intersection over $x \in U$ of the sets
    $\{ g \mid \lVert g(x) \rVert \le 1 \}$; it is closed by
    \texttt{isClosed\_iInter}, since each such set is closed by
    \texttt{isClosed\_le} applied to the continuous maps $g \mapsto \lVert g(x)
    \rVert$ (the evaluation $g \mapsto g(x)$ is continuous by
    \texttt{continuous\_apply}) and the constant map $1$. An intersection of two
    closed sets is closed.
\end{proof}

\begin{lemma}[The image lies in a compact box]
    \label{lem:polar-image-subset-compact-box}
    \uses{def:weak-star-polar, lem:polar-pointwise-bounded}
    Let $U$ be a neighbourhood of $0$ in $E$. By
    Lemma~\ref{lem:polar-pointwise-bounded} choose, for each $x \in E$, a bound
    $r(x) \ge 0$ with $\lVert \varphi(x) \rVert \le r(x)$ for all
    $\varphi \in \operatorname{polar}(U)$. Then the box
    $\prod_{x \in E} [-r(x), r(x)]$ is compact in $E \to \mathbb{R}$, and the
    image $(\cdot)\,''\,\operatorname{polar}(U)$ is contained in it.
\end{lemma}
\begin{proof}
    Each factor $[-r(x), r(x)]$ is a compact interval
    (\texttt{isCompact\_Icc}), so the product box is compact by Tychonoff's
    theorem (\texttt{isCompact\_univ\_pi}). If $g = (\cdot)\,\varphi$ with
    $\varphi \in \operatorname{polar}(U)$, then for each $x \in E$ we have
    $\lVert g(x) \rVert = \lVert \varphi(x) \rVert \le r(x)$, so
    $g(x) \in [-r(x), r(x)]$; hence $g$ lies in the box.
\end{proof}

\begin{lemma}[Compactness of the image]
    \label{lem:polar-image-compact}
    \uses{lem:polar-image-subset-compact-box, lem:polar-image-eq, lem:polar-image-closed}
    Let $U$ be a neighbourhood of $0$ in $E$. The image
    $(\cdot)\,''\,\operatorname{polar}(U)$ is compact in $E \to \mathbb{R}$.
\end{lemma}
\begin{proof}
    By Lemma~\ref{lem:polar-image-subset-compact-box} the image is contained in
    the compact box $\prod_{x \in E} [-r(x), r(x)]$. By
    Lemma~\ref{lem:polar-image-eq} the image equals the intersection
    $L \cap \{ g \mid \forall x \in U,\ \lVert g(x) \rVert \le 1 \}$, which is
    closed by Lemma~\ref{lem:polar-image-closed}. A closed subset of a compact
    set is compact (\texttt{IsCompact.of\_isClosed\_subset}), so the image is
    compact.
\end{proof}

\begin{theorem}[Bourbaki's locally convex Banach--Alaoglu]
    \label{thm:banach-alaoglu-bourbaki}
    \uses{def:weak-star-polar, lem:coe-embedding, lem:polar-image-compact}
    Let $E$ be a real locally convex topological vector space and let $U$ be a
    neighbourhood of $0$ in $E$. Then $\operatorname{polar}(U)$ is compact in
    $\mathrm{WeakDual}\,\mathbb{R}\,E$.
\end{theorem}
\begin{proof}
    \uses{lem:coe-embedding, lem:polar-image-compact}
    By Lemma~\ref{lem:coe-embedding} it suffices to prove that the image
    $(\cdot)\,''\,\operatorname{polar}(U)$ is compact in $E \to \mathbb{R}$,
    which is Lemma~\ref{lem:polar-image-compact}.
\end{proof}
